% =============================================================================
% DCC-KV: Destination-Conditioned Compact KV Communication
% 主稿骨架 — 组装版
%
% 说明：
%   第 1--4 章（摘要、引言、相关工作、问题定义、方法、Algorithm 1）已由用户在
%   桌面 AuthorKit27 (1).pdf 中完成。本文件按 ACM AuthorKit 模板重建其 LaTeX 骨架，
%   使全文可编译；第 5--8 章为本轮续写内容，位于 sections/ 目录。
%
%   第 1--4 章现以本重建版定稿（2026-09-16 决定）：原件即上述桌面 PDF，仓库内不存在
%   原始 .tex，故不再等待替换。公式编号已与原件逐号对齐为 (1)--(27)，引用一律走
%   \eqref；详见 docs/writing_scope_and_metrics.md。
%
% 编译：xelatex main.tex && bibtex main && xelatex main.tex && xelatex main.tex
%      （含中文，必须用 xelatex/lualatex + ctex）
% =============================================================================

\documentclass[sigconf,nonacm]{acmart}

% --- 中文支持 ---------------------------------------------------------------
% 注意：acmart 与 ctex 存在 \baselinestretch 冲突，故直接用 xeCJK。
% 字体为本机已有字体（Windows 简体中文环境）。若在其他平台编译，
% 请把下面的字体名替换为该系统已安装的 CJK 字体。
\usepackage{xeCJK}
\setCJKmainfont{SimSun}[BoldFont=SimHei]
\setCJKsansfont{Microsoft YaHei}
\setCJKmonofont{FangSong}

\usepackage{amsmath,amssymb,amsfonts}
\usepackage{graphicx}
\usepackage{booktabs}
\usepackage{algorithm}
\usepackage{algpseudocode}
\usepackage{listings}

\graphicspath{{figures/}}

% --- 自定义命令：与 PDF 前半部分保持一致 -------------------------------------
\newcommand{\Attn}{\mathrm{Attn}}
\newcommand{\TopK}{\mathrm{TopK}}
\newcommand{\softmax}{\mathrm{softmax}}
\newcommand{\Er}{\mathcal{E}_r}
\newcommand{\Br}{\mathcal{B}_r}
\newcommand{\Es}{\mathcal{E}_s}
\newcommand{\Ds}{\mathcal{D}}
\newcommand{\Cm}{\mathcal{C}}
\newcommand{\eps}{\varepsilon}

% =============================================================================
\begin{document}

\title{DCC-KV：面向超长文本协同推理的目的端条件化紧凑 KV 通信}

% 作者信息与仓库链接暂不写入正文（2026-09-16 决定）。下列字段是**有意的占位符**，
% 由作者本人在投稿前填入；在此之前不要代填。
\author{作者姓名}
\affiliation{%
  \institution{作者单位}
  \country{China}}
\email{author@example.com}

\begin{abstract}
超长文本推理通常通过序列并行将上下文切分至多个设备，但精确全局注意力要求在设备间循环或聚合
完整的 Key--Value（KV）块，使通信量和远程注意力计算随上下文长度线性增长。已有 KV 缓存压缩
方法主要在单设备内依据局部或历史注意力保留重要条目，难以直接回答一个分布式场景中的关键问题：
同一源端上下文应当向具有不同 Query 分布的目的设备发送什么信息？本文提出 DCC-KV
（Destination-Conditioned Compact KV Communication），一种面向长上下文协同推理的训练无关
框架。每个目的设备首先从本地 Query 中选取少量代表样本，并通过 Ring 通信传播其注意力需求；
随后，每个源设备针对不同目的设备分别构造方向非对称的紧凑 KV。受 MIT 团队 Attention Matching
\cite{zweiger2026attentionmatching} 的启发，我们采用 Key 选择、注意力质量校准与 Value 回归
构造紧凑表示。本文的核心创新在于面向变长紧凑 KV 设计异步 All-to-Allv 执行机制：远程块到达后
立即参与注意力计算，并通过满足交换律与结合律的 Online Softmax 状态增量归并，从而使通信与
计算重叠，且计算结果不依赖消息到达顺序。相较于传输完整 KV，DCC-KV 将每条通信边的数据量从
$O(L_s(d_h+d_v))$ 降至 $O(B_{s,r}(d_h+d_v))$，其中 $B_{s,r} \ll L_s$。本文给出完整算法、
误差分析及面向长文本检索、问答与多跳推理的实验方案，为带宽受限的多 GPU 或异构设备协同推理
提供一种查询感知的近似注意力路径。
\end{abstract}

\begin{CCSXML}
<ccs2012>
   <concept><concept_id>10003752.10003809</concept_id>
     <concept_desc>Theory of computation~Distributed algorithms</concept_desc>
     <concept_significance>500</concept_significance></concept>
   <concept><concept_id>10010147.10010178</concept_id>
     <concept_desc>Computing methodologies~Natural language processing</concept_desc>
     <concept_significance>500</concept_significance></concept>
</ccs2012>
\end{CCSXML}

\ccsdesc[500]{Theory of computation~Distributed algorithms}
\ccsdesc[500]{Computing methodologies~Natural language processing}

\keywords{超长文本推理；协同推理；序列并行；KV 缓存压缩；分布式注意力}

\maketitle

% =============================================================================
\section{引言}
\label{sec:intro}

长文档问答、代码仓库分析和多轮智能体推理要求语言模型处理数十万乃至更长的上下文。尽管
FlashAttention \cite{dao2022flashattention,dao2023flashattention2} 等工作降低了单设备注意力的显存访存开销，完整
上下文仍可能超过单卡容量。序列并行由此成为扩展上下文长度的重要路径：输入沿序列维度被划分到
多个设备，每个设备保留局部 Query、Key 和 Value，并通过跨设备通信恢复全局注意力。

Ring Attention \cite{liu2023ring} 通过循环传递 KV 块计算精确注意力，并将通信与块级计算
重叠；DistFlashAttn \cite{li2024distflash}、Striped Attention \cite{brandon2023striped}
和 USP \cite{fang2024usp} 进一步改进了负载均衡或并行拓扑。这些方法避免了近似误差，但每个
设备最终仍需处理所有相关远程 KV。设源设备 $s$ 持有长度为 $L_s$ 的序列块，则其向一个目的
设备发送完整 KV 的数据量为 $O(L_s(d_h + d_v))$。当上下文增长或互联带宽有限时，KV 传输和
远程注意力逐渐成为推理瓶颈。

另一条研究路线通过 KV 选择、驱逐或合并减少缓存规模。H2O \cite{zhang2023h2o}、
Scissorhands \cite{liu2023scissorhands}、SnapKV \cite{li2024snapkv} 与
PyramidKV \cite{cai2024pyramidkv} 均利用注意力稀疏性保留关键状态；Quest \cite{tang2024quest}
则强调关键 KV 与当前 Query 有关。更近期的 Attention Matching
\cite{zweiger2026attentionmatching} 不仅选择 Key，还显式匹配压缩前后的注意力质量与输出，
从而将上下文压缩为可参与后续注意力的潜在表示。然而，现有压缩方法主要面向单设备缓存管理，
通常为所有后续 Query 构造一个共享压缩结果。分布式协同推理中的目的设备持有不同序列块，其
Query 分布也不相同；同一个源端 KV 块对不同目的端的重要部分因而未必一致。

本文研究如下问题：在不交换完整 KV 的前提下，如何让每个设备获得与其本地 Query 相适配的远程
上下文，并保持不同上下文块在全局 Softmax 中的正确竞争关系？我们的核心观察是，协同推理中的
通信不应只压缩源端拥有什么，还应显式编码目的端需要什么。据此，我们提出 DCC-KV。目的设备先
发送少量代表 Query 作为需求摘要；源设备据此为每条有向通信边构造专属紧凑 KV，并分别保持该源
块的注意力质量和语义输出。最终，各块通过可结合的 Online Softmax 状态进行全局归一化。

本文的工作内容与原创边界如下：
\begin{itemize}
\item 在协同推理场景中采用目的端条件化的 KV 通信抽象：对每条有向边 $s \to r$，源设备依据
      目的设备的代表 Query 构造专属紧凑缓存，使 $C^{s \to r_1}$ 与 $C^{s \to r_2}$ 可以不同。
\item 紧凑 KV 构造受 MIT 团队 Attention Matching \cite{zweiger2026attentionmatching} 启发，
      沿用其 Key 选择、质量偏置和 Value 回归思想；\textbf{该压缩机制不作为本文的原创贡献}。
\item 本文原创设计变长紧凑 KV 的异步 All-to-Allv 通信实现，使已到达远程块的注意力计算与
      其余消息传输并行，并利用 Online Softmax 合并算子的交换律与结合律消除对消息接收顺序的
      依赖。
\item 给出全局注意力的块分解、误差界与通信/计算复杂度分析，并设计覆盖长文本检索、多跳推理
      和端到端性能的实验协议。
\end{itemize}

% =============================================================================
\section{相关工作}
\label{sec:related}

\subsection{长上下文分布式注意力}

Ring Attention \cite{liu2023ring} 将序列块分布在环形设备拓扑上，在传递 KV 的同时执行块级
注意力，使可处理长度随设备数扩展。DistFlashAttn \cite{li2024distflash} 针对训练中的通信、
因果负载不均和重计算进行联合优化；Striped Attention \cite{brandon2023striped} 通过条带化
token 分配改善因果注意力的设备利用率；USP \cite{fang2024usp} 统一 Ring 与 Ulysses 风格的
序列并行。上述方法的主要目标是高效实现精确注意力。

推理阶段允许以可控近似换取更低开销。Star Attention \cite{acharya2024star} 通过共享锚块和
块局部注意力降低上下文编码成本；APB \cite{huang2025apb} 在序列并行框架中压缩并传递关键上
下文块。与这些方法相比，DCC-KV 不为所有目的端复用同一远程摘要，而是由每个目的端的 Query
分布决定每条通信边上的紧凑表示；同时显式校准块级注意力质量，使压缩块可与本地完整 KV 在统一
Softmax 下组合。

\subsection{KV 缓存压缩与查询感知稀疏性}

H2O \cite{zhang2023h2o} 根据累计注意力保留 heavy hitters，Scissorhands
\cite{liu2023scissorhands} 利用重要性持续性控制缓存预算，SnapKV \cite{li2024snapkv} 从提示
末端观察窗口推断后续生成所关注的 KV 位置，PyramidKV \cite{cai2024pyramidkv} 根据层间信息
漏斗分配非均匀缓存预算。这些工作说明，少量 KV 条目往往主导注意力输出，但直接丢弃低分条目会
同时改变块内语义和块间归一化质量。StreamingLLM \cite{xiao2024streamingllm} 则揭示了
attention sink 现象并给出流式场景下的保留策略。

Quest \cite{tang2024quest} 指出 KV 关键性依赖当前 Query，并以页级 Key 统计量估计待加载
页面。DCC-KV 与其共享查询感知动机，但关注跨设备通信：目的端通过代表 Query 向源端暴露需求，
源端再离线于该次注意力计算构造可传输的紧凑表示。

MIT 团队提出的 Attention Matching \cite{zweiger2026attentionmatching} 将 KV 压缩表述为同时
匹配注意力输出和注意力质量的问题，并用 Key 选择、非负拟合及 Value 回归构造紧凑缓存。受其
启发，本文采用相同的基本思想构造目的端条件化紧凑 KV，而\textbf{不将压缩方法本身视为原创
贡献}。本文的主要区别在于通信执行层：针对有向设备边上的变长消息设计异步 All-to-Allv 调度，
并以顺序无关的 Online Softmax 状态归并实现通信--计算重叠。

\subsection{协同推理}

Petals \cite{borzunov2023petals} 展示了利用分散设备协同执行大模型推理与微调的可行性，其
重点是跨参与方分配模型层、路由请求并容忍节点动态性。本文研究互补的序列维度协作：设备运行
相同层的不同上下文块，协作对象不是模型层间激活，而是同层注意力所需的远程 KV。该设定尤其
适合单节点多 GPU、机架内多节点或带宽受限的边缘集群。

% =============================================================================
\section{问题定义}
\label{sec:problem}

\subsection{分布式上下文}

给定长度为 $L$ 的序列 $X$，将其划分至 $N$ 个设备：
\begin{equation}
X = [X_0; X_1; \ldots; X_{N-1}], \qquad L = \sum_{i=0}^{N-1} L_i .
\label{eq:split-seq}
\end{equation}
对某一 Transformer 层和 KV Head，设备 $i$ 本地持有
\begin{equation}
Q_i \in \mathbb{R}^{L_i \times d_h},\quad
K_i \in \mathbb{R}^{L_i \times d_h},\quad
V_i \in \mathbb{R}^{L_i \times d_v}.
\label{eq:local-qkv}
\end{equation}
以下推导针对单个 KV Head；多层、多头情况可独立执行并分配不同预算。

令 $\Ds = \{0,\ldots,N-1\}$。目的设备 $r$ 需要关注的远程源集合记为 $\Er$。双向注意力下可取
$\Er = \Ds \setminus \{r\}$；连续分块的因果注意力下，若块索引与序列顺序一致，则
$\Er = \{s \mid s < r\}$。所有有效有向边组成
\begin{equation}
\mathcal{E} = \{(s,r) \mid s \in \Er\}.
\label{eq:edge-set}
\end{equation}

\subsection{块级注意力分解}

对目的设备 $r$ 上的单个 Query $q$，源块 $s$ 的非归一化注意力质量与局部输出分别为
\begin{align}
M_s(q) &= \sum_{j=1}^{L_s} \exp\!\left(\frac{q K_{s,j}^\top}{\sqrt{d_h}}\right),
\label{eq:mass} \\
O_s(q) &= \frac{\sum_{j=1}^{L_s} \exp\!\left(q K_{s,j}^\top / \sqrt{d_h}\right) V_{s,j}}{M_s(q)} .
\label{eq:out}
\end{align}
于是完整全局注意力可写为
\begin{equation}
\Attn(q) = \frac{\sum_{s \in \Br} M_s(q) \, O_s(q)}{\sum_{s \in \Br} M_s(q)} ,
\label{eq:global-attn}
\end{equation}
其中 $\Br = \{r\} \cup \Er$。该分解揭示了压缩远程块必须保留的两个量：块在全局分母中的质量
$M_s$，以及给定进入该块后返回的语义输出 $O_s$。只匹配归一化后的局部输出而忽略质量，会破坏
不同源块之间的相对权重。

% =============================================================================
\section{方法}
\label{sec:method}

\subsection{总体框架}

对每条有效边 $(s,r)$，DCC-KV 构造
\begin{equation}
C^{s \to r} = \left(C_k^{s \to r},\; \beta^{s \to r},\; C_v^{s \to r}\right),
\label{eq:compact}
\end{equation}
其中
\begin{equation}
C_k^{s \to r} \in \mathbb{R}^{B_{s,r} \times d_h},\quad
\beta^{s \to r} \in \mathbb{R}^{B_{s,r}},\quad
C_v^{s \to r} \in \mathbb{R}^{B_{s,r} \times d_v},
\label{eq:compact-shapes}
\end{equation}
且 $B_{s,r} \ll L_s$。$\beta$ 是加入紧凑 Key Logit 的质量补偿偏置。由于构造过程以目的端
Query 为条件，通常有
\begin{equation}
C^{s \to r_1} \neq C^{s \to r_2}.
\label{eq:per-edge-diff}
\end{equation}

框架包含三个逻辑步骤：目的设备提取并交换代表 Query；源设备针对每条边构造紧凑 KV；设备通过
异步 All-to-Allv 交换变长消息，并在消息到达时增量计算和归并注意力状态。

\begin{figure}[t]
\centering
\includegraphics[width=\linewidth]{fig1_architecture.pdf}
\caption{DCC-KV 的总体架构。目的设备仅发送少量代表 Query $\widehat{Q}_r$（虚线，经 Ring
交换）；每个源设备针对每条有向边 $(s,r)$ 独立构造紧凑 KV $C^{s\to r}$（实线）。由于构造以
目的端 Query 为条件，同一源块对不同目的端产生不同表示
（$C^{s\to r_1} \neq C^{s\to r_2}$）。}
\label{fig:arch}
\end{figure}

\subsection{代表 Query 选择}

直接交换 $Q_r \in \mathbb{R}^{L_r \times d_h}$ 会产生随本地序列长度 $L_r$ 线性增长的通信量。
设备 $r$ 因而仅发送 $M \ll L_r$ 个真实 Query。为降低样本选择开销，所有设备由共享随机种子
生成 Rademacher 矩阵 $P \in \mathbb{R}^{d_p \times d_h}$（$d_p \ll d_h$），并在归一化投影
空间中计算距离：
\begin{equation}
P_{a,b} \sim \frac{1}{\sqrt{d_p}}\{\pm 1\},\quad
z_{r,u} = \frac{P q_{r,u}}{\|P q_{r,u}\|_2 + \eps},\quad
D(z_{r,u}, z_{r,v}) = 1 - z_{r,u}^\top z_{r,v}.
\label{eq:proj}
\end{equation}
共享生成使 $P$ 无需传输；投影表示仅用于近似 Query 间的方向差异。

我们采用最远点采样覆盖 Query 分布。以最新 Query 为初始锚点
$\mathcal{S}_r^{(1)} = \{L_r\}$，第 $m$ 轮执行
\begin{equation}
\Delta_{r,u}^{(m)} = \min_{v \in \mathcal{S}_r^{(m)}} D(z_{r,u}, z_{r,v}),\quad
u_{m+1} = \arg\max_{u \notin \mathcal{S}_r^{(m)}} \Delta_{r,u}^{(m)},\quad
\mathcal{S}_r^{(m+1)} = \mathcal{S}_r^{(m)} \cup \{u_{m+1}\},
\label{eq:fps}
\end{equation}
直至 $|\mathcal{S}_r^{Q}| = M$。最终取回原始高维表示
\begin{equation}
\widehat{Q}_r = Q_r[\mathcal{S}_r^{Q}, :] \in \mathbb{R}^{M \times d_h},
\label{eq:repr-q}
\end{equation}
而非传输投影向量。选择复杂度为 $O(L_r d_h d_p + M L_r d_p)$。

所有 $\widehat{Q}_r$ 形状相同，适合使用 $N-1$ 轮 Ring All-Gather。第 $p$ 轮设备 $i$ 向
$[i+1]_N$ 转发当前 Query 摘要，并从 $[i-1]_N$ 接收另一摘要。通信后，源设备 $s$ 获得所有满足
$(s,r) \in \mathcal{E}$ 的 $\widehat{Q}_r$。

\subsection{目的端条件化紧凑 KV}

本节的紧凑 KV 构造受 MIT 团队 Attention Matching \cite{zweiger2026attentionmatching} 启发。
我们将其质量匹配与输出匹配思想用于每条有向设备边 $(s,r)$，以目的端代表 Query 为条件生成通信
消息；\textbf{Key 选择、质量偏置拟合与紧凑 Value 回归并非本文独立提出的压缩算法}。

\subsubsection{紧凑 Key 选择}

源设备 $s$ 以目的端代表 Query 计算块内注意力，并用跨 Query 的 RMS 响应对 Key 排序：
\begin{align}
A^{s \to r} &= \softmax\!\left(\frac{\widehat{Q}_r K_s^\top}{\sqrt{d_h}}\right) \in \mathbb{R}^{M \times L_s},\notag\\
u_j^{s \to r} &= \sqrt{\frac{1}{M}\sum_{a=1}^{M}\left(A_{a,j}^{s \to r}\right)^2},\notag\\
\mathcal{S}_K^{s \to r} &= \TopK\!\left(\{u_j^{s \to r}\}_{j=1}^{L_s},\; B_{s,r}\right),
\label{eq:key-select}
\end{align}
\begin{equation}
C_k^{s \to r} = K_s\!\left[\mathcal{S}_K^{s \to r}, :\right].
\label{eq:ck}
\end{equation}
RMS 聚合保留了仅对少量代表 Query 具有高响应的 Key。

\subsubsection{注意力质量偏置拟合}

仅保留 Key 子集会减小远程块对全局 Softmax 分母的贡献。对代表 Query $\widehat{q}_{r,a}$，定义
完整块的 Logit 与非归一化质量，以及紧凑 Key 的响应矩阵：
\begin{equation}
\ell_{a,k}^{s \to r} = \frac{\widehat{q}_{r,a} K_{s,k}^\top}{\sqrt{d_h}},\qquad
m_a^{s \to r} = \sum_{k=1}^{L_s} \exp\!\left(\ell_{a,k}^{s \to r}\right),\qquad
G_{a,j}^{s \to r} = \exp\!\left(\frac{\widehat{q}_{r,a} (C_k^{s \to r})_{j}^\top}{\sqrt{d_h}}\right),
\label{eq:mass-target}
\end{equation}
其中 $m^{s \to r} \in \mathbb{R}^{M}$，$G^{s \to r} \in \mathbb{R}^{M \times B_{s,r}}$。令
$w_j = \exp(\beta_j) \ge 0$，则带偏置紧凑块的质量为 $\widehat{m} = Gw$。我们通过非负岭回归
拟合
\begin{equation}
w^* = \arg\min_{e^{-3} \le w \le e^{3}} \left\|G^{s \to r} w - m^{s \to r}\right\|_2^2
      + \lambda_\beta \left\|w - \mathbf{1}\right\|_2^2,
\label{eq:beta-fit}
\end{equation}
其中 $\beta_j^{s \to r} = \log(\max\{w_j^*, \eps\})$。
$w_j$ 表示紧凑 Key 所承载的等效注意力质量，可聚合若干被删除 Key 的指数响应；正则项在拟合
信息不足时将其收缩至无偏置状态 $w = \mathbf{1}$。非负约束保证 $w$ 可映射为 Logit 偏置。
$w$ 的箱约束等价于 $\beta_j^{s \to r} \in [-3, 3]$，即单个紧凑 Key 的 Logit 偏置
上下界；该取值沿自 Attention Matching 附录~C.2 对 RMS 选键分支的配置。
实现中应以逐行最大值平移计算指数项，并在拟合目标中保留相应的行缩放，以避免数值溢出。

\subsubsection{紧凑 Value 拟合}

质量匹配只恢复远程块的 Softmax 权重，不能补偿删除 Value 引起的内容误差。我们在代表 Query
上分别计算完整块输出和带偏置紧凑块的归一化权重：
\begin{equation}
Y^{s \to r} = \softmax\!\left(\frac{\widehat{Q}_r K_s^\top}{\sqrt{d_h}}\right) V_s \in \mathbb{R}^{M \times d_v},
\label{eq:Y}
\end{equation}
\begin{equation}
X^{s \to r} = \softmax\!\left(\frac{\widehat{Q}_r (C_k^{s \to r})^\top}{\sqrt{d_h}} + \mathbf{1}_M (\beta^{s \to r})^\top\right) \in \mathbb{R}^{M \times B_{s,r}}.
\label{eq:X}
\end{equation}
紧凑 Value 由多输出岭回归求得：
\begin{equation}
\begin{split}
(C_v^{s \to r})^* &= \arg\min_{C_v} \left\|X^{s \to r} C_v - Y^{s \to r}\right\|_F^2
                   + \lambda_v \|C_v\|_F^2 \\
                  &= \left(X^\top X + \lambda_v I\right)^{-1} X^\top Y .
\end{split}
\label{eq:value-reg}
\end{equation}
因此 $C_v^{s \to r}$ 一般不等于所选 Key 对应的原始 Value；每个紧凑 Value 可共同吸收多个被
删除 Value 在代表 Query 上的输出贡献。质量拟合与 Value 拟合分别逼近块级分母和条件输出。对
任一代表 Query，原始块的未归一化输出可分解为
\begin{equation}
\underbrace{\sum_{k=1}^{L_s} e^{\ell_{a,k}^{s \to r}} V_{s,k}}_{\text{全局 Softmax 分子贡献}}
= \underbrace{m_a^{s \to r}}_{\text{块级质量}}
  \underbrace{\left(\sum_{k=1}^{L_s} \frac{e^{\ell_{a,k}^{s \to r}}}{m_a^{s \to r}} V_{s,k}\right)}_{\text{条件输出}},
\label{eq:decomp}
\end{equation}
这说明同时拟合 $m$ 与 $Y$ 可分别恢复远程块对全局注意力分母和分子的贡献。

\subsection{异步 All-to-Allv 通信与顺序无关的 Online Softmax}

不同有向边采用不同预算 $B_{s,r}$，因此 $C^{s \to r}$ 的长度随源--目的设备对变化。本文采用
循环错位的 All-to-Allv 调度：在第 $p \in \{1,\ldots,N-1\}$ 轮，设备 $s$ 向 $r = [s+p]_N$
发送 $C^{s \to r}$，并从 $[s-p]_N$ 接收消息；若对应边无效，则将发送长度置为零。实现时先
交换每条边的长度与缓冲区偏移，再为各轮提交非阻塞收发。与等待全部通信完成后统一计算不同，
我们使用独立的通信流与计算流，并以逐块完成事件建立依赖：任一远程块到达后，计算流立即启动该
块的注意力计算，通信流则继续接收其余消息，由此形成
\begin{equation}
\text{接收 } C^{s_1 \to r} \;\parallel\; \text{计算 } C^{s_0 \to r} \;\parallel\; \text{发送 } C^{r \to t}
\label{eq:pipeline}
\end{equation}
的流水。\textbf{该实现是本文区别于既有紧凑 KV 方法的核心系统设计。}

\begin{figure}[t]
\centering
\includegraphics[width=\linewidth]{fig2_async_pipeline.pdf}
\caption{异步 All-to-Allv 流水。上：通信流与计算流相互独立，远程块到达后立即启动该块的
注意力计算，无需等待全部消息。下：与"等全部消息到达再统一计算"的同步实现对比；理想流水下
总耗时由两阶段中较慢者支配，而非二者之和。}
\label{fig:async}
\end{figure}

消息的异步到达要求注意力归并与接收顺序无关。对任一块 $b$，令其 Logit 与 Value 为
$\{g_{b,j}, v_{b,j}\}_{j=1}^{T_b}$，其中本地块和远程块分别使用
\begin{equation}
g_{r,j} = \frac{q K_{r,j}^\top}{\sqrt{d_h}},\qquad
g_{s \to r,j} = \frac{q (C_k^{s \to r})_{j}^\top}{\sqrt{d_h}} + \beta_j^{s \to r}.
\label{eq:logits}
\end{equation}
每个块独立生成 Online Softmax 状态
\begin{equation}
\begin{split}
S_b &= (m_b, l_b, o_b),\quad m_b = \max_j g_{b,j},\quad l_b = \sum_j e^{g_{b,j} - m_b},\\
o_b &= \sum_j e^{g_{b,j} - m_b} v_{b,j}.
\end{split}
\label{eq:state}
\end{equation}
对两个状态 $S_a$ 与 $S_b$，沿用在线 Softmax 的归并算子 $\oplus$ \cite{milakov2018online}：
\begin{align}
m_{ab} &= \max(m_a, m_b),\notag\\
l_{ab} &= e^{m_a - m_{ab}} l_a + e^{m_b - m_{ab}} l_b,
\label{eq:merge-params}\\
o_{ab} &= e^{m_a - m_{ab}} o_a + e^{m_b - m_{ab}} o_b,\notag\\
S_a \oplus S_b &= (m_{ab}, l_{ab}, o_{ab}).
\label{eq:merge-op}
\end{align}
该算子等价于对两个块的非归一化 Softmax 分子和分母求和，因而满足
\begin{equation}
S_a \oplus S_b = S_b \oplus S_a,\qquad
(S_a \oplus S_b) \oplus S_c = S_a \oplus (S_b \oplus S_c).
\label{eq:comm-assoc}
\end{equation}
因此，目的设备可按远程消息的实际到达顺序增量更新
\begin{equation}
S_r^{\text{global}} = S_r \oplus \bigoplus_{s \in \Er} S^{s \to r},\qquad
\widehat{\Attn}_r(q) = \frac{o_r^{\text{global}}}{l_r^{\text{global}}} .
\label{eq:global-state}
\end{equation}
交换律保证结果与块的接收顺序无关，结合律允许任意分组或树形归并；只要所有有效块恰好处理一次，
异步执行与等待全部消息后统一计算在数学上等价。与此同时，本地完整 KV 与远程紧凑 KV 始终共享
同一全局 Softmax 分母。

\begin{algorithm}[t]
\caption{DCC-KV 的单层单 KV Head 流程}
\label{alg:dcckv}
\begin{algorithmic}[1]
\Require 本地 $\{Q_i, K_i, V_i\}_{i=0}^{N-1}$，代表数 $M$，预算 $\{B_{s,r}\}$
\Ensure 每个设备 $r$ 的协同注意力输出 $O_r$
\ForAll{设备 $r$（并行）}
  \State 投影 $Q_r$，以最远点采样得到 $\widehat{Q}_r$
\EndFor
\State 通过 Ring 交换所有有效目的端的 $\widehat{Q}_r$
\ForAll{有效边 $(s,r) \in \mathcal{E}$（并行）}
  \State 用式~\eqref{eq:ck} 选取 $C_k^{s \to r}$
  \State 用式~\eqref{eq:beta-fit} 拟合 $\beta^{s \to r}$
  \State 用式~\eqref{eq:value-reg} 重构 $C_v^{s \to r}$
  \State 打包 $C^{s \to r} = (C_k^{s \to r}, \beta^{s \to r}, C_v^{s \to r})$
\EndFor
\ForAll{目的设备 $r$（并行）}
  \State 计算本地状态 $S_r$，并提交变长消息的非阻塞 All-to-Allv 收发
  \ForAll{已到达的 $C^{s \to r}$（按任意顺序）}
    \State 计算 $S^{s \to r}$，并用式~\eqref{eq:merge-op} 立即归并
  \EndFor
  \State 所有接收完成后，用式~\eqref{eq:global-state} 输出 $O_r$
\EndFor
\end{algorithmic}
\end{algorithm}

% =============================================================================
% 以下为本文续写的章节（第 5--8 章）
% =============================================================================
\section{复杂度与误差分析}
\label{sec:analysis}

本节给出 DCC-KV 的通信量与计算复杂度，并将端到端近似误差分解为可独立界定的若干项。分析沿
用第~\ref{sec:problem} 节的记号：$N$ 为设备数，$L_s$ 为源块长度，$d_h$、$d_v$ 分别为
KV Head 的 Key 与 Value 维度，$M$ 为代表 Query 数，$d_p$ 为随机投影维度，$B_{s,r}$ 为边
$(s,r)$ 的压缩预算。所有结论针对单个 Transformer 层、单个 KV Head；多层多头情形按层与头
线性叠加。

\subsection{通信量分析}
\label{sec:analysis-comm}

\subsubsection{消息长度}

DCC-KV 的通信由两段组成。第一段是代表 Query 的环形交换：每轮转发一个
$\widehat{Q}_r \in \mathbb{R}^{M \times d_h}$，进行 $N-1$ 轮，故单设备发送量为
\begin{equation}
\label{eq:comm-query}
V_{\text{query}} = (N-1) \, M \, d_h .
\end{equation}
第二段是变长紧凑 KV 的 All-to-Allv：对每条有效边 $(s,r)$，消息长度为
$B_{s,r}(d_h + d_v) + B_{s,r}$，其中最后一项对应质量偏置 $\beta^{s \to r}$。若以
$|\mathcal{E}_s|$ 记源设备 $s$ 的出边数，则其单层发送总量为
\begin{equation}
\label{eq:comm-kv}
V_{\text{kv}} = \sum_{r \in \mathcal{E}_s}
B_{s,r}\,(d_h + d_v + 1).
\end{equation}

作为对照，精确序列并行下源设备 $s$ 必须向同一组目的端发送完整 KV 块，发送量为
\begin{equation}
\label{eq:comm-dense}
V_{\text{dense}} = |\mathcal{E}_s| \, L_s \, (d_h + d_v) .
\end{equation}

若各边预算相同，记为 $B$，则当 $B \ll L_s$ 且 $N$ 不太大时有
\begin{equation}
\label{eq:comm-ratio}
\frac{V_{\text{query}} + V_{\text{kv}}}{V_{\text{dense}}}
\approx \frac{B}{L_s} + \frac{(N-1) M d_h}{|\mathcal{E}_s| L_s (d_h+d_v)} .
\end{equation}
即压缩比主要由 $B/L_s$ 决定，代表 Query 的开销随设备数线性增长但与上下文长度无关，
在长上下文区间可忽略。

\begin{figure}[t]
\centering
\includegraphics[width=\linewidth]{fig4_comm_scaling.pdf}
\caption{单边通信量的解析对比（$d_h+d_v=256$，$B=64$）。精确序列并行的通信量随源块长度
$L_s$ 线性增长；DCC-KV 的紧凑 KV 传输量 $O(B_{s,r}(d_h+d_v))$ 与 $L_s$ 无关，二者相差
$L_s/B$ 倍。代表 Query 的摘要开销 $O(NMd_h)$ 随设备数增长，但与上下文长度无关。
\emph{本图为解析结果，非实测数据。}}
\label{fig:comm}
\end{figure}

\subsubsection{与集合通信实现的匹配}

式~\eqref{eq:comm-query} 的形状规整（所有 $\widehat{Q}_r$ 同形），可直接由 $N-1$ 轮
Ring All-Gather 实现。式~\eqref{eq:comm-kv} 的关键性质是各边长度不同，标准 All-to-All
要求每个接收端预知各来源的消息长度。4.4 节的两阶段实现（先交换长度与缓冲区偏移，再提交
非阻塞收发）把这一约束转化为一次 $O(N^2)$ 的元数据交换，该交换与消息体传输解耦，因此不
进入主通信路径。

\subsection{计算复杂度分析}
\label{sec:analysis-compute}

\subsubsection{代表 Query 选择}

按式~\eqref{eq:proj}--\eqref{eq:repr-q}：随机投影为 $O(L_r d_h d_p)$；$M$ 轮最远点采样，每轮维护并更新当前子集
到所有候选点的最小距离，代价为 $O(M L_r d_p)$。故选择开销为
\begin{equation}
\label{eq:cost-fps}
C_{\text{select}} = O\!\left(L_r d_h d_p + M L_r d_p\right),
\end{equation}
与 $L_r$ 成线性、与 $B_{s,r}$ 无关。注意该开销在所有设备上并行发生。

\subsubsection{紧凑 KV 构造}

对每个目的端 $r$，源设备 $s$ 需计算响应矩阵
$A^{s\to r} = \mathrm{softmax}(\widehat{Q}_r K_s^\top / \sqrt{d_h}) \in
\mathbb{R}^{M \times L_s}$，代价 $O(M L_s d_h)$；RMS 聚合与 Top-$B$ 选择为
$O(M L_s)$ 与 $O(L_s \log B)$。质量偏置拟合的求解规模为 $M \times B$，
投影梯度下降的每次迭代代价 $O(M B)$；Value 回归为 $M \times B$ 的岭回归，
闭式解的构造成本 $O(M B^2 + M B d_v)$。合计
\begin{equation}
\label{eq:cost-build}
C_{\text{build}} = O\!\left(M L_s d_h + M B_{s,r}(d_h + d_v) + M B_{s,r}^2\right).
\end{equation}
式~\eqref{eq:cost-build} 的第一项是主导项。这是目的端条件化的直接代价：源设备 $s$ 需为其
每个出边重复一次构造，总构造代价乘以 $|\mathcal{E}_s|$。

\subsubsection{目的端注意力计算}

目的端 $r$ 的本地块注意力代价为 $O(L_r^2 d_h)$（因果情形减半）。对每个到达的远程紧凑块，
计算式~\eqref{eq:logits} 的 Logit 与式~\eqref{eq:state} 的状态，代价为 $O(L_r B_{s,r}(d_h + d_v))$。全部远程块合计
\begin{equation}
\label{eq:cost-attn}
C_{\text{attn}} = O\!\left(L_r^2 d_h + L_r (d_h + d_v) \textstyle\sum_{s \in \mathcal{E}_r} B_{s,r}\right).
\end{equation}
与精确远程注意力 $O(L_r \sum_{s} L_s (d_h+d_v))$ 相比，式~\eqref{eq:cost-attn} 的第二项
关于 $\sum_s B_{s,r}$ 而非 $\sum_s L_s$，这是压缩在计算侧的收益。

\subsubsection{总览与权衡}

表~\ref{tab:complexity} 汇总各阶段复杂度。需要强调的是，DCC-KV 并不降低本地块的注意力
代价，也不减少层数或头数；其收益集中在跨设备通信量与远程注意力计算量两处，同时以
$|\mathcal{E}_s|$ 倍的构造代价与 $M$ 倍的代表 Query 通信作为代价。当 $B \ll L_s$
且 $M \ll L_r$ 时，这一交换在长上下文区间是划算的。

\begin{table}[t]
\centering
\small
\caption{DCC-KV 各阶段的单层单头复杂度。$B \equiv B_{s,r}$，$L \equiv L_s \approx L_r$。}
\label{tab:complexity}
\begin{tabular}{lll}
\toprule
阶段 & DCC-KV & 精确序列并行 \\
\midrule
代表 Query 交换 & $O(N M d_h)$ & --- \\
紧凑 KV 构造 & $O(|\mathcal{E}_s|(M L d_h + M B^2))$ & --- \\
跨设备通信量 & $O(|\mathcal{E}_s| B (d_h+d_v))$ & $O(|\mathcal{E}_s| L (d_h+d_v))$ \\
远程注意力 & $O(L \sum_s B)$ & $O(L \sum_s L_s)$ \\
\bottomrule
\end{tabular}
\end{table}

\subsection{误差分解与误差界}
\label{sec:analysis-error}

设目的端 $r$ 上某一 Query $q$ 的精确全局注意力为 $\mathrm{Attn}(q)$，DCC-KV 的输出为
$\widehat{\mathrm{Attn}}(q)$。按式~\eqref{eq:global-attn} 的块分解，误差来自三个可分离的环节。以下均针对
单个 Query，且假设所有有效块恰好被处理一次。

\subsubsection{误差的三个来源}

\paragraph{(i) Query 摘要误差。}
源设备仅看到 $\widehat{Q}_r$ 而非完整 $Q_r$。若 $q$ 在投影空间中与所选代表 Query 距离
较远，则据此构造的紧凑 KV 未必适配 $q$。该误差由代表 Query 集合对 $Q_r$ 分布的覆盖度
决定，记作 $\epsilon_{\text{repr}}$。

\paragraph{(ii) 块级质量匹配误差。}
式~\eqref{eq:beta-fit} 在非负约束与 $\lambda_\beta$ 正则下拟合 $m_a^{s\to r}$，一般不能精确复现。定义
\begin{equation}
\label{eq:err-mass}
\epsilon_{\text{mass}}(q) = \sum_{s \in \mathcal{E}_r}
\left| \widehat{M}_s(q) - M_s(q) \right| ,
\end{equation}
其中 $\widehat{M}_s(q) = \sum_{j=1}^{B_{s,r}} \exp(g_{s\to r,j})$。该误差直接改变远程块在
全局 Softmax 分母中的相对权重，是所有误差项中对最终输出影响最直接的一项。

\paragraph{(iii) 条件输出误差。}
式~\eqref{eq:value-reg} 以 $\lambda_v$ 岭回归拟合 $Y^{s \to r}$，得到的是代表 Query 上的最小二乘解，
不保证对任意 $q$ 成立。定义
\begin{equation}
\label{eq:err-out}
\epsilon_{\text{out}}(q) = \sum_{s \in \mathcal{E}_r}
M_s(q) \left\| \widehat{O}_s(q) - O_s(q) \right\|_2 .
\end{equation}

\subsubsection{误差界}

将式~\eqref{eq:global-attn} 的分子分母分别扰动。设
$\widehat{M}_s = M_s + \delta M_s$、$\widehat{O}_s = O_s + \delta O_s$，且记
$S(q) = \sum_{s \in \mathcal{B}_r} M_s(q)$。当 $\sum_s |\delta M_s| < S/2$ 时，
\begin{equation}
\label{eq:err-bound}
\left\| \widehat{\mathrm{Attn}}(q) - \mathrm{Attn}(q) \right\|_2
\;\le\;
\frac{\epsilon_{\text{out}}(q)}{S(q)}
\;+\;
\frac{2\, \epsilon_{\text{mass}}(q)}{S(q)} \cdot \max_{s \in \mathcal{B}_r} \|O_s(q)\|_2 .
\end{equation}
式~\eqref{eq:err-bound} 表明：质量误差通过"相对权重被放大"的机制进入输出，且放大系数是
块输出范数的量级；输出误差则以与自身质量成比例的权重进入。二者均以全局 Softmax 分母
$S(q)$ 归一化，因此远程块占比越大，单块误差的绝对影响越小。

\textbf{该界的首次数值检验。}式~\eqref{eq:err-bound} 此前只被陈述、从未被检验，
而 \S\ref{sec:analysis-error} 的性质 3 又用它论证 DCC-KV 在设备数较大时仍可用。
E13 在「一个压缩源块 $+$ 一个精确上下文块」的两块归并设定下逐 Query 计算两侧：
共 $5760$ 个 Query（$5$ 个种子 $\times$ 关注强度 $\{0,8\}$ $\times$
$B\in\{8,\dots,256\}$ $\times$ $\beta\in\{\text{none},\text{full}\}$），
\textbf{零个反例}，紧致度 $\mathrm{RHS}/\mathrm{LHS}\in[1.00,3.51]$、
中位 $2.10$。判据带 $10^{-9}$ 的相对容差，用于吸收两条数值路径的舍入：
当 $\delta M\equiv0$ 时该界取\textbf{等号}
（$\mathrm{RHS}=\epsilon_{\text{out}}/S=\mathrm{LHS}$），但左侧经逐块 online 归并、
右侧由解析式算出，故纯浮点意义上的「$<$」共 $189$ 例、全部落在取等号的
$10$ 个格子里。若不加该容差，就会把「界恰好取等号」误报成「界被推翻」。

前置条件 $\sum_s|\delta M_s|<S/2$ 的违反率则\textbf{不可忽略}，且其分布本身是结论：
$\beta$ 全量时最高 $16.7\%$，而\emph{不加} $\beta$ 时最高 $100\%$，
并沿预算单调下降（$B=8$ 为 $91.0\%$，$16$ 为 $70.6\%$，$32$ 为 $48.5\%$，
$64$ 为 $15.2\%$，$B\ge128$ 为 $0$）。即\emph{界在误差最大的区间
（小预算且无质量补偿）本就没有保证}，这与该区间的误差主要由 $\beta$ 的缺失
主导完全一致。$\beta$ 的作用因此不止于降低 $\epsilon_{\text{mass}}$：
它同时把该界的适用范围从小预算侧扩展开来。

\begin{figure}[t]
\centering
\includegraphics[width=\linewidth]{fig3_error_decomposition.pdf}
\caption{端到端误差的来源与误差界。质量匹配误差 $\epsilon_{\text{mass}}$ 与输出匹配误差
$\epsilon_{\text{out}}$ 分别作用于块级分母与条件输出，二者均以全局 Softmax 分母 $S(q)$
归一化，故不随设备数发散。代表 Query 覆盖误差 $\epsilon_{\text{repr}}$ 不显式出现在界中，
而是通过前两项间接体现。}
\label{fig:err}
\end{figure}

式~\eqref{eq:err-bound} 中不含 $\epsilon_{\text{repr}}$ 的显式项，因为 Query 摘要误差
通过 $\epsilon_{\text{mass}}$ 与 $\epsilon_{\text{out}}$ 间接体现：$q$ 偏离代表 Query
越远，源设备为它构造的 $(\widehat{M}_s, \widehat{O}_s)$ 就越不准确。这一耦合意味着
$M$ 的选取需要在"摘要开销小"与"$\epsilon_{\text{mass}}, \epsilon_{\text{out}}$ 低"
之间权衡，且该权衡是可测量的（见第~\ref{sec:experiment} 节 E2/E5）。

\subsubsection{若干可陈述的性质}

\paragraph{性质 1（质量守恒的充分条件）。}
若对某条边 $(s,r)$ 有 $B_{s,r} = L_s$，则式~\eqref{eq:key-select} 的 Top-$B$ 选择保留全部 Key，式~\eqref{eq:ck}
退化为 $C_k^{s \to r} = K_s$；此时式~\eqref{eq:beta-fit} 取 $w = \mathbf{1}$、$\beta = \mathbf{0}$ 是
零残差解，式~\eqref{eq:value-reg} 取 $C_v^{s \to r} = V_s$ 亦为零残差解。故 DCC-KV 在预算不裁剪时精确
退化为全注意力，即该框架是精确注意力的一个严格泛化。

需要强调该命题是\textbf{存在性}的充分条件，而非「按默认拟合结果实现即精确」：
预算不裁剪时式~\eqref{eq:value-reg} 的岭回归解 $\widehat{C}_v$ 一般\emph{不等于}
$V_s$，输出侧仍有余项。E2 在 $B=L_s=256$ 处实测：质量侧相对误差为 $0$
（不加 $\beta$）与 $5\times10^{-3}$（$\beta$ 全量），而输出侧相对误差为
$0.070$--$0.225$。也就是说\emph{预算不裁剪时主导误差来自 Value 通路}，
而非选键或 $\beta$——这与 E2 在 $B<L_s$ 处的误差分解定性一致，
也说明「精确退化」只对质量侧无条件成立。

\paragraph{性质 2（无偏置下的质量下界）。}
由于 $w \ge 0$ 且式~\eqref{eq:key-select} 按 RMS 响应选择 Key，被保留 Key 的响应在代表 Query 上系统性地
偏高。故在无偏置情形（$\beta = \mathbf{0}$）下通常有 $\widehat{M}_s(q) \le M_s(q)$，
即压缩后远程块的全局权重被系统性低估，而非随机波动。这为 4.3.2 节引入质量补偿偏置提供
了直接的动机：$\beta$ 需要做的正是把被裁掉的指数响应重新分配给保留的 Key。

\paragraph{性质 3（误差不随块数累积发散）。}
式~\eqref{eq:err-bound} 的右侧以 $S(q)$ 归一化，而 $S(q)$ 本身随有效块数增长。因此，即使
每条边都引入非零误差，只要各边的质量误差不同向叠加，总输出误差不会随设备数 $N$ 线性
放大。这一点是 DCC-KV 在 $N$ 较大时仍可用的必要条件。需要说明的是，该性质依赖误差的
符号分布；若所有边系统性低估质量（如性质 2 所述），则需靠 $\beta$ 拟合抵消，这也解释了
4.3.2 节中正则项选择 $w = \mathbf{1}$ 而非 $w = \mathbf{0}$ 作为收缩目标的设计理由。

\textbf{机制归因的一处更正。}上文把「不随块数放大」归给"各边误差不同向叠加"。
E13 在「每块长度固定、块数 $N$ 增长」的设定下实测（每块 $64$ token、
每块预算 $16$，故总长随 $N$ 增长而每块的压缩强度不变）：
$N$ 由 $1$ 增至 $8$ 时输出侧偏差\emph{不增}（$\beta$ 全量、关注强度 $8$ 下
$0.208\to0.128$；不加 $\beta$、强度 $0$ 下 $0.322\to0.360$，仅 $1.12$ 倍），
而同期 $\sum_s|\delta M_s|/S$ 单调\emph{上升}（$0.038\to0.078$ 与 $0.250\to0.580$），
且 $|\sum_s\delta M_s|\,/\sum_s|\delta M_s|$ 在不加 $\beta$ 时\textbf{恒为} $1.000$——
即各块误差\textbf{完全同向叠加}，一次抵消都没有发生。故真正阻止发散的是
\textbf{分母 $S(q)$ 随有效块数增长}这一项；符号抵消只是在此之上的一层\emph{附加}
裕度（$\beta$ 全量、$N=8$ 时该比值降至 $0.576$），是充分条件而非必要条件。
上文把符号抵消写成条件的表述因此\emph{过强}：即便各边误差同向叠加，
只要 $S(q)$ 同步增长，结论依然成立。

\subsection{异步流水的理论加速上限}
\label{sec:analysis-async}

设单条边的紧凑 KV 传输耗时为 $T_{\text{comm}}$，其对应的注意力计算耗时为
$T_{\text{comp}}$。在纯串行实现（等全部消息到达后统一计算）下，$N-1$ 轮的总耗时约为
\begin{equation}
\label{eq:serial-time}
T_{\text{serial}} \approx (N-1)\left( T_{\text{comm}} + T_{\text{comp}} \right).
\end{equation}
在理想流水下（通信流与计算流完全独立，且任意时刻均有可计算的已到达块），总耗时受两阶段
中较慢者支配：
\begin{equation}
\label{eq:pipeline-time}
T_{\text{async}} \approx \max \Big( (N-1) T_{\text{comm}},\;(N-1) T_{\text{comp}} \Big)
+ T_{\text{comp}} .
\end{equation}
因此理想加速比上界为
\begin{equation}
\label{eq:speedup-bound}
\frac{T_{\text{serial}}}{T_{\text{async}}}
\;\lesssim\;
\frac{T_{\text{comm}} + T_{\text{comp}}}{\max(T_{\text{comm}}, T_{\text{comp}})}
\;=\; 1 + \frac{\min(T_{\text{comm}}, T_{\text{comp}})}{\max(T_{\text{comm}}, T_{\text{comp}})} .
\end{equation}
当 $T_{\text{comm}} \approx T_{\text{comp}}$ 时上界趋近 $2\times$；当两者严重失衡时加速
趋于 $1$。该结论给出一个重要的实践含义：DCC-KV 的异步设计在通信与计算预算相当的区间收益
最大，而在任一侧主导的区间收益有限。第~\ref{sec:experiment} 节 E6 的加速比测量应显式报告
$T_{\text{comm}} / T_{\text{comp}}$ 的实测值，以便判断实测加速偏离上界的原因是流水不充分
还是预算本身失衡。

需要指出，式~\eqref{eq:speedup-bound} 是理论上界，不含通信启动延迟、元数据交换开销、
负载不均衡以及调度抖动。此外，式~\eqref{eq:comm-assoc} 保证的只是数值结果与接收顺序无关，而在线 Softmax
状态归并本身引入浮点舍入差异；该差异的量级已由 E0 界定（\S\ref{sec:exp-cpu}）：
$\mathrm{FP32}$ 下不超过 $7.9\times10^{-7}$、$\mathrm{FP64}$ 下处于机器舍入量级，
且不随乱序试验次数增长；低精度（$\mathrm{FP16}$/$\mathrm{BF16}$）下的失效边界仍需 E8。
因此，异步实现的
实测加速比应低于式~\eqref{eq:speedup-bound}，且二者之差可用于衡量实现质量。

\subsection{小结}

本节确立三点。第一，DCC-KV 把每条边的通信量从 $O(L_s(d_h+d_v))$ 降至
$O(B_{s,r}(d_h+d_v))$，代价是 $|\mathcal{E}_s|$ 倍的构造开销与 $O(NM d_h)$ 的摘要开销，
二者在长上下文区间均可接受。第二，误差可分解为质量误差与输出误差两项并以全局 Softmax
分母归一化，因此不随设备数发散；质量误差通过相对权重的放大机制进入输出。第三，异步流水
的理想加速上界由 $T_{\text{comm}}$ 与 $T_{\text{comp}}$ 的比值决定，该比值是可测的，因而
加速比与上界的差距可以归因。

需要明确的是，本节给出的误差界依赖 $\lambda_\beta$、$\lambda_v$ 与代表 Query 覆盖度等
量。其\emph{紧致程度}已在合成数据上被测量（见上文 E13 与第~\ref{sec:discussion} 节）：
$\mathrm{RHS}/\mathrm{LHS}$ 中位 $2.10$、$95$ 分位 $3.51$、$5760$ 个 Query 零反例，
故它是量级可用的界；但该比值、以及 $\epsilon_{\text{mass}}$ 与 $\epsilon_{\text{out}}$ 的
经验分布，均由合成分布决定，\textbf{在真实模型与真实长文本上尚未测量}。
第~\ref{sec:experiment} 节给出测量方案，其中可在纯 CPU 环境
下复现的部分已标注，需要 GPU 的部分明确标为待执行。

\section{实验方案}
\label{sec:experiment}

本节给出 DCC-KV 的实验设计、评测协议与结果报告规范。需要首先明确本节的性质：\textbf{本文
尚未报告端到端任务指标与通信性能。}截止撰写时，仓库所实现并已在 CPU 上验证的部分为
（i）紧凑 KV 构造全流程，（ii）在线 Softmax 归并算子的顺序无关性，（iii）同步版分布式
注意力的数值等价性，（iv）变长 All-to-All 消息的打包、解包与接口形状。异步通信流水与多卡
可扩展性依赖 GPU 环境，尚未执行。因此，本节按"已可复现的验证"与"待执行的实验"两类分别
陈述，并给出后者的完整协议，以便结果补齐后可直接填入。

\subsection{实验设置}
\label{sec:exp-setup}

\subsubsection{硬件与软件}

CPU 验证部分在任意多核环境即可复现，无需 GPU。GPU 部分按 \texttt{docs/reproducibility.md}
第 4 节设定的配置执行：M2（多卡同步）最低 2$\times$A100 80G，M3（异步）最低
4$\times$A100 80G NVLink，M4（端到端 8B）推荐 8$\times$A100 80G，70B/72B 与跨节点实验
需要 8$\times$H100 或双节点环境。所有依赖版本通过 \texttt{requirements.lock} 锁定；
PyTorch 版本、CUDA、NCCL 与 Transformers 版本均记入 \texttt{ExperimentMetadata}
（见 \S\ref{sec:exp-metadata}）。

\subsubsection{模型}

评测覆盖四个模型以检验方法对模型族的迁移性：
Llama-3.1-8B-Instruct、Llama-3.1-70B-Instruct、Qwen2.5-72B-Instruct、
Mistral-7B-Instruct-v0.3。前两者为同族不同规模，用于检验规模迁移；Qwen2.5 与 Mistral 为
不同训练语料与架构配置，用于检验跨族迁移。

\subsubsection{数据集}

长上下文任务覆盖检索、问答、多跳推理与长文困惑度四类：
LongBench 与 LongBench-v2（多任务长文理解）、$\infty$Bench（超长文本）、
RULER（合成检索与多跳，可控长度）、Needle-in-a-Haystack（单针检索，长度可控）、
PG-19 与 wikitext-103（语言建模困惑度）。其中 RULER 与 NIAH 的长度可精确控制，
是绘制"误差--上下文长度"曲线的首选。

\subsubsection{基线}

三个基线：Ring Attention（精确序列并行的代表，作为精度上界与性能参照）、
FastKV（共享压缩，即所有目的端复用同一压缩结果，用于隔离"目的端条件化"的贡献）、
APB（全网共享 anchor 的压缩方案）。三者与 DCC-KV 的差异在通信语义上是清晰的：
Ring Attention 不压缩，FastKV 压缩但不条件化，APB 压缩且共享锚点，
DCC-KV 压缩且逐边条件化。

\subsubsection{超参数与随机性}

关键超参数为代表 Query 数 $M$、投影维度 $d_p$、压缩预算 $B_{s,r}$、
正则系数 $\lambda_\beta$ 与 $\lambda_v$。所有实验固定随机种子（默认 42，备选 123 / 1024 / 7），
每个数据点至少重复 10 次，报告 median 与 p5/p95，并给出 bootstrap 95\% 置信区间。

\subsection{评测协议与元数据规范}
\label{sec:exp-metadata}

为避免"实验条件不可追溯"导致的结论不可复现，每次 run 由 \texttt{ExperimentMetadata}
记录：\texttt{run\_id}、\texttt{git\_commit}、\texttt{pytorch\_version}、
\texttt{model\_name}、\texttt{context\_length}、\texttt{gpu\_count}、\texttt{seed}、
\texttt{budget\_ratio}、\texttt{sync\_async}、\texttt{num\_repr\_queries}、
\texttt{projection\_dim}、\texttt{lambda\_beta}、\texttt{lambda\_value}、
\texttt{rope\_extension\_used} 与 \texttt{rope\_extension\_disclosed}。

其中后两项的设计值得说明：若为达到目标上下文长度而使用了 RoPE 外推，必须显式披露。这一约束
的目的是防止把"位置编码外推带来的效果"误归因为压缩机制的贡献。凡 \texttt{rope\_extension\_used}
为真而 \texttt{rope\_extension\_disclosed} 为假的 run，其结果不进入主表。

结果聚合由 \texttt{RunResult.from\_values} 完成，统一产出 median、p5/p95 与
bootstrap 95\% CI。该实现意味着所有报告数字必须来自多次 run 的分布，而非单次运行值。

\subsection{已验证的 CPU 实验}
\label{sec:exp-cpu}

以下四组实验对应仓库中已存在的测试，均可在纯 CPU 环境复现。

\subsubsection{E0：顺序无关性}

\textbf{目的}：验证式~(26) 的交换律与结合律在有限精度下成立，即异步乱序归并与顺序归并的
结果差异处于浮点舍入量级。

\textbf{方法}：构造 $K$ 个随机块（长度与 $d_v$ 各异），各自生成在线 Softmax 状态，然后
（i）按原顺序归并，（ii）随机置换后归并，（iii）按树形分组归并，比较三者输出。

\textbf{已有证据}：\texttt{verify\_order\_invariance} 在 $K=8$、20 次随机试验下返回真；
\texttt{test\_order\_invariance} 与 \texttt{test\_merge\_two\_states} 通过，两状态归并与
单次 dense 计算的差异在 \texttt{FP64} 下低于 $10^{-6}$。

\textbf{待补}：需把"置换数量"与"归并树形状"作为自变量，报告最大相对误差随二者的变化曲线，
以界定异步实现的数值安全边界。

\subsubsection{E1：紧凑 KV 构造的接口与形状不变量}

\textbf{目的}：确认 $C_k^{s\to r}$、$\beta^{s\to r}$、$C_v^{s\to r}$ 的形状与预算约束一致，
且随机构造可复现。

\textbf{已有证据}：\texttt{build\_compact\_kv} 在 $B=64$、$d_h=d_v=32$ 下返回
\texttt{keys} $(64,32)$、\texttt{logit\_bias} $(64,)$、\texttt{values} $(64,32)$、
\texttt{selected\_indices} $(64,)$；Top-$B$ 选择返回的索引互不重复；
相同种子下代表 Query 选择完全一致。

\subsubsection{E2：压缩保真度}

\textbf{目的}：测量 $\epsilon_{\text{mass}}$ 与 $\epsilon_{\text{out}}$ 在受控压缩比下的
经验分布。

\textbf{方法}：固定 $L=256$、$d_h=d_v=32$、$\mathrm{FP64}$，扫描
$B \in \{8, 16, 32, 64, 128\}$，对每个 $B$ 计算紧凑块与完整块在代表 Query 上的
质量误差与输出误差，并与 Dense 全注意力对比。

\textbf{已有证据}：\texttt{test\_fast\_kv\_vs\_dense\_within\_tolerance} 在 $B=64$ 下
测得压缩方法与 Dense 的最大绝对偏差 $< 0.1$；基线 APB 在更激进压缩下偏差 $< 0.3$。
这两个数字给出了压缩误差在随机张量设定下的量级参照。

\textbf{待补}：现有测试只断言"小于阈值"，未报告误差随 $B$ 的变化趋势。需要补充的是
$\epsilon_{\text{mass}}(B)$ 与 $\epsilon_{\text{out}}(B)$ 的完整曲线，用于验证
\S\ref{sec:analysis-error} 中性质 2 的预测（无偏置时质量被系统性低估）。

\subsubsection{E3：边级条件化 vs 共享压缩}

\textbf{目的}：为 H1 与 H2 提供数值侧证。

\textbf{方法}：在相同预算下比较 DCC-KV（逐边条件化）与 FastKV（共享压缩）相对 Dense 的
误差。H1 断言同一源块对不同目的端的紧凑 KV 显著不同，以 KL 散度 $> 0.5$ 为判据。

\textbf{已有证据}：\texttt{test\_dcc\_kv\_lower\_error\_than\_shared} 在 $B=32$ 的强压缩
设定下同时计算两种方法的误差。需要如实说明：该测试的当前断言是对二者的误差\emph{分别}
设上界（均 $< 0.5$），\textbf{并未断言 DCC-KV 严格优于 FastKV}。测试注释中亦注明
"CPU mock 不一定严格"。因此这一项目前只构成接口与量级验证，不构成 H2 的证据。

\textbf{待补}：需要（i）把断言从"各自有界"提升为"配对比较且 DCC-KV 显著更优"，
（ii）把 $B$ 与目的端 Query 分布差异作为自变量扫描，（iii）把 KL 散度量化为 H1 的直接证据。
这是与本文核心主张最相关、也最需要补强的一环。

\subsubsection{E4：分布式等价性}

\textbf{目的}：验证同步版 DCC-KV 与完整分布式注意力在数值上一致，且变长消息在真实多进程
环境下正确传递。

\textbf{已有证据}：Ring Attention 与 Dense 的最大偏差 $< 10^{-5}$（$L=256$，四块，
$\mathrm{FP64}$）；同步版 DCC-KV 与 Dense 的偏差断言阈值为 $10^{-2}$；
两进程 gloo 后端下 \texttt{all\_reduce}、\texttt{broadcast} 与变长
\texttt{all\_to\_all\_v} 的数值结果均逐元素正确。

\textbf{说明}：$10^{-2}$ 与 $10^{-5}$ 的差异反映了压缩引入的近似误差，而非实现错误。
该阈值本身即是一个可报告的结论：在 $L=256$、$B=64$ 的设定下，同步 DCC-KV 相对精确注意力的
最大绝对偏差在 $10^{-2}$ 量级。

\subsection{待执行的 GPU 实验}
\label{sec:exp-gpu}

以下实验依赖 GPU 与多卡环境，尚未执行。协议如下，结果补齐后填入对应表格。

\subsubsection{E5：消融实验}

消融设计覆盖组件与预算两个维度。

\paragraph{A1 通信集大小。}扫描有效边数 $|\mathcal{E}_s|$，观察构造开销与通信开销随设备数
的分配比例，用于检验式~\eqref{eq:cost-build} 中 $|\mathcal{E}_s|$ 倍构造代价的实际影响。

\paragraph{A2 压缩预算。}在 5 档预算下测量任务指标与通信量，绘制精度--通信量帕累托前沿。
该曲线是本文最核心的可视化结果：它同时呈现压缩收益与精度代价，并与 FastKV、APB 的同类
曲线叠加比较。

\paragraph{A3 组件拆分。}四个变体：完整 DCC-KV、（i）移除质量偏置 $\beta$、
（ii）移除 Value 回归（直接取所选 Key 对应的原始 Value）、（iii）二者同时移除
（退化为纯 Key 选择）。按 H3，移除 $\beta$ 应导致 $\ge 0.5$ 点退化，移除 Value 回归应导致
$\ge 1.0$ 点退化。

\paragraph{A5 异步 vs 同步。}固定其他条件，比较异步流水与同步实现的 p50 延迟。按 H4，
异步应带来 $\ge 1.05\times$ 的 p50 加速。此外应同时报告
$T_{\text{comm}} / T_{\text{comp}}$ 的实测值，以便按式~\eqref{eq:speedup-bound} 判断
实测加速比与理论差距的归因。

\subsubsection{E6：主表与可扩展性}

主表为 $3$ 模型 $\times$ $5$ 上下文长度 $\times$ $2$ GPU 数 $\times$ $2$ 同步模式
$\times$ $3$ 基线，共 $144$ 个数据点，每点 $\ge 10$ 次 run。报告指标包括任务精度、
prefill 延迟（p50 与 p95）、端到端吞吐与跨设备通信量。

可扩展性按 H5 检验：设备数翻倍时加速应 $\ge 1.5\times$。若未达标，需按
\S\ref{sec:exp-negative} 的规范转为负结果讨论。

\subsubsection{E7：负结果与适用边界}
\label{sec:exp-negative}

按仓库发布清单的要求，以下四种条件必须显式报告结果，无论其是否符合预期：

\emph{短上下文（$L < 4$K）。}此时 $B/L$ 不再小，压缩的固定开销（代表 Query 交换与逐边
构造）占比上升，预期 DCC-KV 无收益甚至劣于精确注意力。这一边界需明确给出，而非回避。

\emph{低预算（$B < 1\%$）。}此时 $\epsilon_{\text{mass}}$ 与 $\epsilon_{\text{out}}$ 均
显著增大，用于确定方法的可用压缩下界。

\emph{强检索任务。}检索任务对精确的 token 级定位敏感，压缩可能系统性地丢失关键条目。
这是对"误差不随块数发散"这一性质的压力测试。

\emph{batch = 1。}此时计算侧难以饱和，$T_{\text{comp}}$ 偏小，
式~\eqref{eq:speedup-bound} 预测异步收益受限。该场景直接检验流水的收益边界。

\subsubsection{E8：数值稳定性}

把 \S\ref{sec:exp-cpu} 的 E0 扩展到 $\mathrm{FP16}$/$\mathrm{BF16}$、$K \ge 32$ 个块与
更深的归并树结构，测量顺序无关性在低精度下的失效边界。该实验的结果决定异步实现能否在
真实推理精度下安全使用，是工程可行性的关键一环。

\subsection{结果报告规范}

本节确立三条报告规范，用于约束后续结果的呈现方式。

\textbf{第一，区分主张来源。}凡涉及紧凑 KV 构造机制（Key 选择、质量偏置、Value 回归）的
实验结果，应在叙述中明确标注该机制源自 Attention Matching，DCC-KV 的贡献在于把它条件化到
有向边上。涉及异步通信与顺序无关归并的结果才是本文原创主张的证据。

\textbf{第二，区分已执行与待执行。}如表~\ref{tab:exp-status}，每项实验的状态必须显式标注。
凡属"待执行"的项目，正文中不得出现未标注的具体数值。

\textbf{第三，负结果与正结果同等报告。}按发布清单，若 H2 / H3 / H4 / H5 中任一未达标，
主张即按预设阈值收敛，并在本节如实报告。这一机制的作用是防止在结果不利时事后调整主张范围。

\begin{table}[t]
\centering
\small
\caption{实验状态总览。"CPU" 表示可在纯 CPU 环境复现，"GPU" 表示依赖多卡环境。}
\label{tab:exp-status}
\begin{tabular}{llll}
\toprule
编号 & 内容 & 环境 & 状态 \\
\midrule
E0 & 顺序无关性（$\mathrm{FP64}$） & CPU & 已有证据 \\
E1 & 构造接口与形状不变量 & CPU & 已有证据 \\
E2 & 压缩保真度（单点） & CPU & 部分证据 \\
E3 & 边级条件化 vs 共享压缩 & CPU & 仅量级验证 \\
E4 & 分布式等价性 & CPU（gloo） & 已有证据 \\
E5 & 消融 A1/A2/A3/A5 & GPU & 待执行 \\
E6 & 主表与可扩展性 & GPU & 待执行 \\
E7 & 负结果与适用边界 & GPU & 待执行 \\
E8 & 低精度数值稳定性 & GPU & 待执行 \\
\bottomrule
\end{tabular}
\end{table}

\subsection{本节小结}

本节的实验设计有两个特点值得指出。第一，它把"可验证"与"待验证"严格分离：E0--E4 给出的是
接口正确性、数值等价性与顺序无关性的证据，这些证据不依赖 GPU，因而是当前状态下唯一可对外
声明的结论；E5--E8 给出的是完整协议，其结果决定了本文关于精度与性能的主张能否成立。
第二，它把与本文核心主张最相关的检验（E3）明确标为"仅量级验证"——当前测试并未断言
DCC-KV 优于共享压缩。这一诚实标注是必要的：若把"两者误差都有界"读作"边级条件化更优"，
就是把尚未测得的结论当作已有证据。

因此，本节当前能支撑的结论限于：紧凑 KV 构造流程在实现层面自洽且可复现；在线 Softmax
归并算子在 $\mathrm{FP64}$ 下满足顺序无关性；变长 All-to-All 消息在多进程环境下正确传递；
同步版分布式注意力与 Dense 的偏差处于压缩误差量级。关于任务质量、通信收益、多卡可扩展性
以及与基线的优劣对比，均待 GPU 实验完成后报告。

\section{讨论}
\label{sec:discussion}

\subsection{本文主张的边界}

在讨论本文的意义之前，需要先把主张的范围划定清楚，因为 DCC-KV 涉及的三个环节来源不同。

\textbf{不属于本文贡献的部分。}将 KV 缓存压缩为可参与后续注意力的紧凑表示，并同时匹配
注意力质量与输出，这一思想来自 MIT 团队的 Attention Matching。本文 4.3 节的 Key 选择、
质量偏置拟合与 Value 回归三步，在其思想框架内实现，\textbf{本文不主张该压缩机制为原创}。
同样，块级注意力分解、在线 Softmax 的 $(m,l,o)$ 状态表示与归并算子，均是既有工作
（FlashAttention、Online Normalizer Calculation）的组成要素。本文把后者从"单设备内分块的
数值技巧"重新理解为"跨设备异步消息的顺序无关归并机制"，但算子本身不是新的。

\textbf{属于本文贡献的部分。}其一是目的端条件化的通信抽象：对每条有向边 $(s,r)$ 单独构造
$C^{s\to r}$，使同一源块对不同目的端的紧凑表示可以不同，且显式校准块级注意力质量以维持
跨块在全局 Softmax 中的竞争关系。其二是面向变长紧凑 KV 的异步 All-to-Allv 执行机制：
环形错位调度、独立通信流与计算流、以逐块到达事件建立依赖，实现通信与计算的重叠。
其三是在此设定下给出的通信/计算复杂度分析与误差分解。这三项中，第二项是唯一在既有文献中
没有直接对应物的系统设计。

这一划分意味着本文的可证伪点也是明确的：若 E5/A2 显示逐边条件化相对共享压缩没有带来
可测的精度收益，则第一项贡献失去意义；若 E5/A5 显示异步流水相对同步实现没有带来延迟收益，
则第二项贡献失去意义。二者不相互担保。

\subsection{为什么需要目的端条件化}

一个自然的替代方案是让所有目的端共享同一份压缩 KV。这种做法更简单：源设备只需构造一次，
通信量也更低（不需要按边区分）。DCC-KV 放弃了这一优惠，换取的是每边独立的表示。

这一交换是否划算，取决于一个可测量的事实：同一源块对不同目的端的注意力分布究竟差异多大。
若差异很小，逐边构造的额外开销就是浪费；若差异显著，共享压缩就会系统性地保留对某些目的端
无用、而对另一些目的端关键的条目。H1（KL $> 0.5$）正是为这一事实设定的判据，而 E3 是唯一
能回答它的实验。

值得指出的是，这一问题的答案与工作负载强相关。在因果分块设定下，靠后的目的端与靠前的
目的端所关注的远程内容本就不同；在多跳推理任务中，不同位置需要的证据片段也不同。相反，
在需要全局平均信息的摘要类任务中，各目的端的注意力分布可能高度相似，此时条件化的收益会
明显下降。因此本文不主张目的端条件化在所有任务上都有收益，而是主张其收益的上界由 Query
分布差异决定，且这一差异可测。

\subsection{异步收益的条件与代价}
\label{sec:discussion-tension}

\S\ref{sec:analysis-async} 给出的加速上界 $1 + \min(T_{\text{comm}}, T_{\text{comp}}) /
\max(T_{\text{comm}}, T_{\text{comp}})$ 说明了一件事：异步不是无条件有益的。当通信显著慢于
计算时，加速上界趋近 $1$；当计算显著慢于通信时同样趋近 $1$。只有在两者相当时，上界才
接近 $2\times$。

这一结论对本文的叙事有一个不受欢迎但必须承认的推论。\textbf{压缩本身会减小
$T_{\text{comm}}$，因而可能把系统推离"通信与计算相当"的区间，反而降低异步的边际收益。}
也就是说，压缩与异步这两个设计之间存在张力：压缩越激进，异步能够重叠掉的时间越少。
DCC-KV 的真实收益上限来自"压缩后的通信 + 重叠"这一组合，而非二者各自收益的简单相加。
在设计实验时若只单独报告压缩比与异步加速比，会掩盖这一交互。第~\ref{sec:experiment} 节的消融 A4 正是为测定该交互而设：它同时扫预算与同步模式，并据此判定 A2 与 A5 的边缘结果能否分开引用。

异步还引入若干工程代价。变小长度消息使通信模式不规则，标准 All-to-All 需要预知接收长度，
因此须增加一次元数据交换；乱序到达要求归并算子满足交换律与结合律，这在数学上由式~\eqref{eq:comm-assoc} 保证，
但在有限精度下只近似成立（E8 的目标）；独立的通信流与计算流占用额外资源，在计算侧本已饱和
的场景下可能相互争抢。这些代价不体现在复杂度表达式中，但会体现在实测数据里。

\subsection{误差界的意义与局限}

式~\eqref{eq:err-bound} 的价值在于把端到端误差拆成质量误差与输出误差两项，并显示二者均以
全局 Softmax 分母 $S(q)$ 归一化。这带来一个有实践意义的推论：远程块占比越大、误差被吸收得
越好，因此长上下文场景比短上下文更适合该方法。这也解释了为什么 E7 预期短上下文是负结果区间。

但这一误差界有三点局限。第一，它是\emph{条件性}的：推导假设
$\sum_s |\delta M_s| < S/2$。E13 已实测该条件的违反率，它不是一个小概率事件：
$\beta$ 全量时最高 $16.7\%$、不加 $\beta$ 时最高 $100\%$，
且沿预算单调下降（$B=8$ 时 $91.0\%$ 降到 $B\ge128$ 的 $0$）。
故「界成立」与「界有保证」是两件事：界可以在前置条件不满足的格子上依然成立，
但那属于经验观测而非推导结论，二者必须分开陈述。
第二，它的紧致程度\textbf{已在合成数据上被测量}：E13 得到
$\mathrm{RHS}/\mathrm{LHS}$ 中位 $2.10$、$95$ 分位 $3.51$、最小 $1.00$
（$5760$ 个 Query 中零个反例），故它是一个量级可用的界而非装饰；
但该比值由合成分布决定，与真实模型上是否同阶仍未经测量——
$\epsilon_{\text{mass}}$ 与 $\epsilon_{\text{out}}$ 的机制级经验分布
（E2 的完整曲线，以及 E2b / E10 / E11 的网格）同样只到这一步。第三，它把代表 Query 的覆盖误差隐含在
$\epsilon_{\text{mass}}$ 与 $\epsilon_{\text{out}}$ 之中，因此没有给出 $M$ 与误差之间
显式的关系。该关系目前只由实验回答：E9 在 $M$--$B$ 网格上分离出「$B$ 决定可达上界、
$M$ 决定能否触及上界」的职责划分，但这是合成数据上的结论，且 $M$ 的取值随任务与预算
移动，尚不存在与任务无关的理论公式。

\subsection{「质量相近」的判定}
\label{sec:quality-comparable}

\S\ref{sec:exp-gpu} 的 H2 要求「DCC-KV 相对共享压缩的质量提升 $\ge 1.5$ 个百分点」
\emph{且}「在质量相近前提下 prefill 加速 $\ge 1.10\times$」。前半句的比较对象是明确的；
后半句的前提「质量相近」却只是一个词：与谁相近、由谁判定、容差多大，规划文档都没有规定。
本节给出一个可操作的定义。之所以要把它当作一个独立问题，是因为\textbf{它不是一个措辞问题，
而是一个会改变结论的问题}——容差取宽一点，任何实现都能宣称自己「质量相近」。

\paragraph{参照物取精确注意力，而不是共享压缩基线。}最省事的读法是让「质量相近」与前半句
同源，即要求 DCC-KV 与共享压缩基线在质量上相近。这一读法\textbf{自我矛盾}：若两者的质量差
被限制在 $\delta$ 以内，同时又要求二者相差至少 $1.5$ 个百分点，则只有当
$\delta > 1.5$ 个百分点时才可能同时成立；而那时「提升 $1.5$ 个百分点」已经落在「相近」的
范围之内，前半句不再表达任何主张。因此参照物必须\emph{换一个}，取\textbf{精确注意力}
（未经压缩的全局注意力）。换过之后 H2 的三段含义互不重叠、可以同时成立：压缩后的 DCC-KV
与精确注意力\textbf{不劣}（压缩基本无损），DCC-KV 相对共享压缩\textbf{更高}（逐边条件化
确有收益），且 prefill \textbf{更快}（系统设计确有收益）。

\paragraph{判定主体是协议与检验程序，不是作者。}「相近」必须由\textbf{预先固定的留出评测
协议} $\mathcal{B}$ 与一个\textbf{事先声明的检验程序}共同判定，不由作者事后目测，也不按
结果挑选指标。在 $\mathcal{B}$ 下取得 DCC-KV 与精确注意力的任务质量差
$\Delta = Q_{\text{DCC}} - Q_{\text{dense}}$ 及其 $95\%$ 置信区间，判据为
$\mathrm{CI}_{\text{low}}(\Delta) > -\delta$，即\textbf{单侧非劣检验}：压缩允许有代价，
但代价不得超过 $\delta$。采用区间而非点估计，是因为点估计为 $0$ 可能只是样本噪声恰好抵消；
用区间才能把「测不出差异」与「确实无差异」分开。这里取单侧而非双侧等价，理由是该前提的
\emph{功能}是排除「以质量换速度」这一混淆——只要 DCC-KV 没有显著更差，加速比的比较就是
公平的；若它同时显著更优（$\mathrm{CI}_{\text{low}}(\Delta) > 0$），前提以更强的形式成立，
但此时必须在结果中\textbf{显式披露方向}，不得笼统写作「相近」。

\paragraph{容差 $\delta$ 被一个下界与两个上界夹住。}下界来自统计分辨率：$\delta$ 不得小于
\textbf{同一配置重复 run 的噪声底线}（可用 run 间质量差的 p95$-$p50 或 bootstrap 区间
半宽估计）。若 $\delta$ 小于噪声底线，「相近」将永远无法被建立——测不出差异与确实无差异
不可区分，判定退化为对噪声的断言。上界之一是先验可定的：\textbf{非劣边界不得大于所要检测
的效应量}，故必须有 $\delta < 1.5$ 个百分点，否则「相近」宽到把前半句所声称的提升一并
吞没。上界之二是需要数据的：$\delta$ 还必须小于\textbf{共享压缩相对精确注意力的退化量}
$\Delta_{\text{bad}} = Q_{\text{dense}} - Q_{\text{share}}$，否则共享压缩本身也会落入
「相近」的判定范围，该前提就失去了把两者区分开的能力。于是
\begin{equation}
\label{eq:delta-window}
\text{噪声底线} \;\le\; \delta \;<\; \min\{\,1.5,\ \Delta_{\text{bad}}\,\}\ \text{个百分点}.
\end{equation}
在区间内的取值依据是保守性：先保证判定不是噪声伪影（下界），再不让它失去分辨力（上界）。
实现层面，这一区间被写成一条\textbf{硬断言}——等价检验函数在 $\delta$ 越出
\eqref{eq:delta-window} 时\emph{拒绝执行}，而不是返回「不相近」。这个区别是必要的：容差
取值不当意味着\textbf{本次实验不具备判定该前提的分辨率}，与「质量确实不相近」是两个不同的
结论，前者必须报告为分辨率不足，不能被静默读成后者。

\paragraph{两个前置条件必须与加速比测量绑定。}其一，\textbf{配对性}：质量侧与性能侧必须取自
\textbf{同一组 run}（同 batch、同提示、同上下文长度、同种子）。否则「不劣」与「更快」可能
来自不同条件，二者的组合没有意义。其二，\textbf{跨上下文长度的聚合方式}必须先定：E6 在
$4$ 个上下文长度上各产出一组质量差，是取全部长度的联合区间，还是逐长度判定后取交集，会
得到不同结论。本稿采用保守规则——\textbf{四个长度上分别做非劣检验，全部通过才算前提成立}，
任一长度不通过即单独报告该长度，不得用其余长度的通过把它平均掉。

\paragraph{失效模式。}若实测噪声底线达到或超过 $1.5$ 个百分点（在长上下文任务上并非不可能），
\eqref{eq:delta-window} 为空区间，H2 在现有测量精度下\textbf{不可检验}——此时正确的处理是
报告「分辨率不足、无法判定 H2」，而不是判定 H2 不成立。若
$\mathrm{CI}_{\text{low}}(\Delta) \in (-\delta, 0]$，则前提成立而方向未定于「更优」，
此时应如实写作「与精确注意力的差异未超过 $\delta$」，不得写作「与精确注意力相当」——
后者会隐含一个并未被检验的等号。

上述规则已落到代码：\texttt{hypotheses.py} 的 \texttt{h2\_pass\_}\allowbreak\texttt{across\_}\allowbreak\texttt{lengths}
把逐长度判定聚成单一结论，且\textbf{「分辨率不足」与「未达标」分开记账}——
某个长度的前提无法判定时记入 \texttt{unresolved} 而不计入 \texttt{failed}，
整体结果为「不可判定」；\texttt{quality\_comparable\_non\_inferior} 在
$\delta$ 越出可行区间时\emph{抛错而非返回 False}，使上面那条区别成为可执行的。
E6 的结果结构已接上该入口，其两个参数以\emph{无默认值}的命令行选项暴露：
有默认值等于假称前提永远成立。

\subsection{适用场景}

综合以上分析，DCC-KV 的适用条件可以陈述为四项同时成立：
（i）上下文足够长，使 $B/L$ 足够小，压缩带来的通信节省超过其构造与摘要开销；
（ii）跨设备带宽受限，使 $T_{\text{comm}}$ 与 $T_{\text{comp}}$ 处于可比区间，异步有发挥空间；
（iii）源块对不同目的端的注意力分布差异显著，使条件化有实际收益；
（iv）任务对 token 级精确定位的敏感度不高，使压缩误差可被容忍。

四条中任一条不成立，DCC-KV 相对精确序列并行或共享压缩的收益即会消失。这一条件集合比
"面向长上下文协同推理"的泛化表述更窄，但它更可能经得起复现。

反过来说，最不利的场景是可预期的：短上下文、高带宽互联（如单机 NVLink 全连接）、
同构目的端 Query 分布、强检索任务。在这些场景下，精确 Ring Attention 或简单共享压缩
是更合理的选择。

\section{结论}
\label{sec:conclusion}

本文提出了 DCC-KV，一种面向超长文本协同推理的目的端条件化紧凑 KV 通信框架。其基本观察是：
在序列并行设定下，源设备向不同目的设备发送的信息不应相同，因为目的端的 Query 分布不同，
而源块中各条目对不同 Query 分布的重要性也不同。据此，DCC-KV 对每条有向通信边 $(s,r)$
构造专属的紧凑 KV $C^{s\to r} = (C_k^{s\to r}, \beta^{s\to r}, C_v^{s\to r})$，
在压缩的同时显式校准块级注意力质量，并通过满足交换律与结合律的在线 Softmax 状态归并，
使已到达的远程块可以被立即计算，从而实现异步通信与计算的重叠。

本文的主要工作包括四项。第一，给出协同推理场景下的目的端条件化 KV 通信抽象，使同一源块
对于不同目的端产生不同的紧凑表示。第二，设计面向变长紧凑 KV 的异步 All-to-Allv 执行机制，
通过环形错位调度与独立的通信/计算流实现重叠，并以顺序无关的状态归并消除对消息到达顺序的
依赖。第三，给出全局注意力的块级分解，并据此将端到端误差分解为质量误差与输出误差两项，
导出以全局 Softmax 分母归一化的误差上界，说明误差不随设备数发散。第四，给出完整的实验协议，
覆盖长文本检索、问答、多跳推理与困惑度任务，并显式设定负结果报告机制。

在实现层面，紧凑 KV 构造流程、在线 Softmax 归并算子、变长消息传递与同步版分布式注意力
的数值等价性均已有可在纯 CPU 环境复现的验证，具体包括：归并算子在 $\mathrm{FP64}$ 下满足
顺序无关性，与单次密集计算的差异低于 $10^{-6}$；变长 All-to-All 在两进程 gloo 环境下
逐元素正确；\emph{不压缩且不加偏置}的同步版 DCC-KV 相对精确注意力的最大绝对偏差为
$2.5\times10^{-7}$（$\mathrm{FP32}$，$L=256$），说明通信与归并路径本身数值干净，
而压缩设定下的偏差来自紧凑构造的近似。

需要明确说明的是，本文尚未报告端到端任务指标与通信性能。异步通信流水与多卡可扩展性依赖
GPU 环境，尚未执行；关于任务质量、通信收益以及与基线的优劣对比，有待第~\ref{sec:experiment}
节所述实验完成后报告。此外，紧凑 KV 的构造机制沿用 Attention Matching 的 Key 选择、
质量偏置拟合与 Value 回归思想，本文不主张该机制为原创；本文的原创贡献集中于目的端条件化的
通信抽象与变长消息的异步执行机制。

后续工作有三个方向。其一，把代表 Query 的选择从"覆盖 Query 分布"推广到"覆盖对压缩误差
最敏感的方向"，即由误差反馈驱动摘要选择，而非由 Query 几何结构驱动。其二，把逐边独立构造成
本从 $|\mathcal{E}_s|$ 倍降低，例如对 Query 分布相近的目的端做聚类后共享部分构造中间量。
其三，把异步执行机制推广到流水线并行与专家并行的混合场景，检验顺序无关归并在更复杂通信
拓扑下的适用性。

\section*{可复现性说明}

本工作的代码与实验配置见随论文发布的代码仓库；本文对应的 commit hash 与逐次变更
记录见仓库的 \texttt{docs/}\allowbreak\texttt{commit\_}\allowbreak\texttt{log.md}。
依赖版本的冻结清单（\texttt{requirements.lock}）随代码发布提供；每个 run 的完整元数据由
\texttt{ExperimentMetadata} 记录，包括 git commit、PyTorch/CUDA/NCCL 版本、
模型名称、上下文长度、GPU 数量、随机种子、压缩预算比、同步/异步模式、
代表 Query 数与投影维度、两个正则系数，以及 RoPE 外推的使用与披露标志。
所有结果由 \texttt{RunResult} 聚合，报告 median 与 p5/p95 及 bootstrap 95\% 置信区间。
纯 CPU 部分的验证可由 \texttt{pytest tests/} 直接复现；GPU 部分的最低硬件配置见
\texttt{docs/reproducibility.md} 第 4 节。若在实验中使用了位置编码外推技术，
其使用情况会在对应结果中一并披露。


% =============================================================================
\bibliographystyle{ACM-Reference-Format}
\bibliography{refs}
\label{sec:refs}

\end{document}
